\reading{Kam-Almaz v. United States}
\readingcite{682 F.3d 1364 (Fed. Cir. 2012)}

\opinion \textsc{Lourie}, Circuit Judge.\ldots

%Majd Kam-Almaz appeals from the final decision of the United States Court of
%Federal Claims dismissing his breach of contract and Fifth Amendment taking
%claims. We affirm.
%
%
%
%\readinghead{BACKGROUND}
%
%Kam-Almaz alleged the following facts in his pleadings before the Court of
%Federal Claims. \textit{See} Compl., Jan. 5, 2009, ECF No. 1; Am. Compl., Jan.
%25, 2010, ECF No. 23.
Kam-Almaz is a United States citizen employed in the
business of international disaster relief assistance. On April 7, 2006,
Kam-Almaz returned home from an overseas business trip. At Dulles International
Airport in Loudoun County, Virginia, Agent Craig Muldowan of the United States
Immigration and Customs Enforcement (``ICE'') detained Kam-Almaz, informing him
that he was a ``person of interest to ICE.'' Muldowan seized
Kam-Almaz's laptop and two flash drives for review by ICE. Before Muldowan
seized the equipment, however, Kam-Almaz informed him that it contained the only
copies of his business files; in response, Muldowan permitted Kam-Almaz to copy
and retain one computer file. Upon seizing the equipment, Muldowan provided to
Kam-Almaz a signed Customs Form 6051D indicating that the equipment would be
detained for up to thirty days. Muldowan also verbally assured Kam-Almaz that
the equipment would be held for no more than seven days.

While the laptop was detained, its hard drive failed, destroying much of
Kam-Almaz's business software. On May 15, 2006, a representative from the U.S.
Customs and Border Protection (``Customs'') sent Kam-Almaz a letter seeking to
assure him that a prompt resolution of the issue would be addressed by Muldowan.
On June 21, 2006, about ten weeks after its seizure, the laptop was returned to
Kam-Almaz. On June 24, 2006, the Director of Investigations for ICE sent
Kam-Almaz a letter representing that ``ICE has made  every attempt to minimize
the inconvenience to [Kam-Almaz]. ICE copied the files and provided them to Mr.
Kam-Almaz during the initial border stop.''

On January 5, 2009, Kam-Almaz filed suit in the Court of Federal Claims\ldots.
%alleging breach of an implied-in-fact contract. On January 25, 2010, he amended
%his complaint and included a takings claim.
Kam-Almaz alleged damages totaling
\$469,480.00 due to lost business contracts resulting from his inability to
access his computer files as well as replacement hardware, software, and
warranty costs. [The Court of Federal Claims dismissed the complaint, and
Kam-Almaz appealed.]
% On June 30, 2010, the government moved to dismiss, arguing that
%the Court of Federal Claims lacked jurisdiction over Kam-Almaz's complaint, and
%that the complaint failed to state a claim upon which relief could be granted.

%In a decision dated January 7, 2011, the Court of Federal Claims granted the
%government's motion to dismiss. \textit{Kam-Almaz v. United States}, 96 Fed.Cl.
%84, 86 (2011). On the breach of contract claim, the court dismissed under Rule
%12(b)(6) of the United States Court of Federal Claims (``RCFC 12(b)(6)'') for
%failure to state a claim upon which relief can be granted. The court held that
%the complaint failed in several respects to allege facts sufficient to find a
%bailment contract. The court found that, because the complaint stated that his
%property was seized, Kam-Almaz did not ``deliver[] personalty'' to the
%government as a bailment requires. \textit{Id.} at 88. The court further found
%that the complaint failed to allege that the government promised to return the
%computer in accordance with Kam-Almaz's instructions or to guard or carefully
%handle the equipment. Again focusing on the fact that the complaint described
%the government's act as a ``seizure,'' the court found that the complaint
%failed to allege facts demonstrating the required mutuality of intent between
%the parties. Finally, the court held that the complaint lacked the necessary
%allegations of Muldowan's authority to enter into a bailment contract.

%The Court of Federal Claims also dismissed Kam-Almaz's takings claim under RCFC
%12(b)(6). The court explained that property seized and retained pursuant to the
%government's police power is not taken for a public use within the context of
%the Fifth Amendment's Takings Clause. Noting that border agents do not have
%authority to seize property without having reasonable cause to suspect a
%violation of law, the court rejected Kam-Almaz's theory that the laptop was
%seized, not according to the police power, but according to an administrative
%border search for security purposes. The court further explained that, if the
%seizure was unauthorized, then the Court of Federal Claims would lack
%jurisdiction, because due process and Fourth Amendment claims are reserved for
%district courts. Finally, the court held that even assuming the government's
%actions were authorized and that an unreasonable delay in returning the
%property amounted to a taking, the court lacks jurisdiction over damage claims
%for due process violations.
%
%Kam-Almaz appealed. We have jurisdiction pursuant to 28 U.S.C. \S~1295(a)(3).



\readinghead{Discussion}

\ldots.

%To avoid dismissal for failure to state a claim under RCFC 12(b)(6), ``a
%complaint must allege facts `plausibly suggesting (not merely consistent with)'
%a showing of entitlement to relief.'' \textit{Acceptance Ins. Cos., Inc. v.
%United States}, 583 F.3d 849, 853 (Fed.Cir.2009) (quoting \textit{Bell Atl.
%Corp. v. Twombly}, 550 U.S. 544, 557, 127 S.Ct. 1955, 167 L.Ed.2d 929 (2007)).
%The facts as alleged ``must be enough to raise a right to relief above the
%speculative  level, on the assumption that all the allegations in the complaint
%are true (even if doubtful in fact).'' \textit{Twombly}, 550 U.S. at 555, 127
%S.Ct. 1955 (citations omitted). At the same time, a court is ``\,`not bound to
%accept as true a legal conclusion couched as a factual allegation.'\,''
%\textit{Id.} (quoting \textit{Papasan v. Allain}, 478 U.S. 265, 286, 106 S.Ct.
%2932, 92 L.Ed.2d 209 (1986)). We review \textit{de novo} a decision to dismiss a
%complaint for failure to state a claim under RCFC 12(b)(6). \textit{Hearts Bluff
%Game Ranch, Inc. v. United States}, 669 F.3d 1326, 1328 (Fed.Cir.2012). We also
%review \textit{de novo} the grant or denial of a motion to dismiss for lack of
%jurisdiction. \textit{See} \textit{Frazer v. United States}, 288 F.3d 1347, 1351
%(Fed.Cir.2002).
%
%
%
%\readinghead{I}

We first address Kam-Almaz's claim for breach of an implied-in-fact bailment
contract.\ldots
%``An implied-in-fact contract with the government requires proof of
%(1) mutuality of intent, (2) consideration, (3) an unambiguous offer and
%acceptance, and (4) actual authority on the part of the government's
%representative to bind the government in contract.'' \textit{Hanlin v. United
%States}, 316 F.3d 1325, 1328 (Fed.Cir. 2003) (internal quotation marks
%omitted). An implied-in-fact contract is founded upon a meeting of the minds
%and ``\,`is inferred, as a fact, from the conduct of the parties showing, in
%the light of the surrounding circumstances, their tacit understanding.'\,''
%\textit{Id.} (quoting \textit{Balt. \& Ohio R.R. v. United States}, 261 U.S.
%592, 597, 58 Ct.Cl. 709, 43 S.Ct. 425, 67 L.Ed. 816 (1923)).
``A bailment relationship is
said to arise where an owner, while retaining title, delivers personalty to
another for some particular purpose upon an express or implied contract. The
relationship includes a return of the goods to the owner or a subsequent
disposition in accordance with his instructions.'' \textit{Lionberger v. United
States}, 371 F.2d 831, 840 (Ct. Cl. 1967); \textit{see also} 19 Williston on
Contracts \S~53:1 (4th ed. 2012) (defining a bailment as ``a delivery of
personalty for some particular purpose, or on mere deposit, upon a contract,
express or implied, that after the purpose has been fulfilled it shall be
redelivered to the person who delivered it, or otherwise dealt with according to
his directions, or kept until he reclaims it, as the case may be.'').\ldots
% Kam-Almaz, as the plaintiff, bears the
%burden of proving the existence of an implied-in-fact contract. \textit{See}
%\textit{Hanlin v. United States}, 316 F.3d 1325, 1330 (Fed. Cir.2003).

%Kam-Almaz contends that the Court of Federal Claims erred by concluding that
%his complaint failed to allege facts sufficient to assert a plausible claim for
%breach of an implied bailment contract.
Kam-Almaz asserts that his complaint
plausibly alleges a bailment because Kam-Almaz ``delivered personalty'' by
transferring possession of his laptop to ICE, and because his computer was
expected to be held only for a short period of time and then returned in working
condition.\ldots
%Kam-Almaz further contends that his complaint plausibly alleges a
%breach of an implied-in-fact contract because the facts alleged demonstrate (1)
%mutuality of intent to contract, (2) a negotiation demonstrating offer,
%acceptance, and consideration, and (3) actual authority of Muldowan to bind the
%United States in contract. Finally, he contends that Supreme Court and Federal
%Circuit cases recognize that a party in Kam-Almaz's shoes may be able to
%establish a breach of an implied contract when property temporarily detained by
%the government is damaged.
%
%The government, in response, contends that the Court of Federal Claims correctly
%dismissed Kam-Almaz's complaint on the  pleadings. According to the government,
%Kam-Almaz did not allege mutual intent to contract because the government
%lawfully seized the laptop pursuant to its police power. The government further
%asserts that Kam-Almaz failed to establish any other element of an
%implied-in-fact bailment contract, including an offer, acceptance,
%consideration, and actual authority to contract.
%
%We agree with the government that the Court of Federal Claims correctly
%dismissed Kam-Almaz's complaint for failure to allege facts plausibly suggesting
%a breach of an implied-in-fact bailment contract. Like the Court of Federal
%Claims, we find numerous deficiencies in Kam-Almaz's complaint. The complaint
%does not plausibly allege the required elements of a bailment.
[The court disagreed.] Kam-Almaz did not
voluntarily ``deliver'' his equipment to Muldowan. \textit{See} \textit{Goudy \&
Stevens, Inc. v. Cable Marine, Inc.}, 924 F.2d 16, 18 (1st Cir. 1991). Rather,
as Kam-Almaz repeatedly alleged, it was involuntarily ``seized.''
The complaint further fails to
allege facts indicating the mutual intent required for an implied-in-fact
contract. A seizure, essentially by definition, lacks mutual intent. Thus, as
the government correctly points out, a seizure pursuant to the government's
authority to police the border generally will not give rise to an
implied-in-fact bailment contract.
\textit{See} \textit{Llamera v. United
States}, 15 Cl. Ct. 593, 597 (1988)\ldots.
%; \textit{see also} \textit{Alde, S.A. v.
%United States}, 28 Fed.Cl. 26, 30 (1993) (collecting cases and observing that
%``[t]hese cases evince a uniform reluctance to find an implied bailment
%contract ... where plaintiff's property has been seized pursuant to the
%Government's exercise of its police power'').
Further, because Kam-Almaz did not voluntarily
deliver his property to the government, his complaint fails to allege any valid
consideration.
%\textit{See} \textit{Llamera}, 15 Cl.Ct. at 598.
In summary, we
agree with the Court of Federal Claims that ``[t]he `purely unilateral act' of
seizing a person's personal property does not evidence intent to enter into a
bailment contract.''
%\textit{Kam-Almaz}, 96 Fed.Cl. at 88 (quoting
%\textit{Alde}, 28 Fed.Cl. at 31).

In support of his bailment claim, Kam-Almaz points to Muldowan's statement that
the laptop would be seized ``for no more than seven days'' and the document
receipt stating that ``shipments may be detained for up to thirty (30) days.''
Those estimates of when the seized property might be returned to
Kam-Almaz, however, are insufficient to allege a bailment contract.
\textit{See, e.g.}, \textit{Llamera}, 15 Cl. Ct. at 597 (noting that even though
the plaintiff received ``a receipt'' for his seized property, the facts ``did
not evidence any intent to enter into a bailment agreement'').
Furthermore, as the relevant
regulation indicates, the fact that Kam-Almaz received a receipt in no way
diminishes the characterization of Muldowan's actions as a unilateral seizure.
\textit{See} 19 C.F.R. \S~162.21(a) (``A receipt for seized property shall be
given at the time of seizure to the person from whom the property is
seized.'').\ldots
%Kam-Almaz also alleges that correspondence between him and the government
%supports the plausibility of his claim for an implied-in-fact bailment contract.
%But those letters do not evidence the government's intent to enter into an
%enforceable implied-in-fact contract with Kam-Almaz. \textit{See}
%\textit{Hanlin}, 316 F.3d at 1330 n. 3.

We therefore conclude that Kam-Almaz failed to plausibly allege a mutual intent
to contract\ldots
%, as an implied-in-fact contract requires,
and further failed to
plausibly allege a voluntary delivery of property, as needed for a bailment
contract.\ldots
% We need not address the other deficiencies in Kam-Almaz's  pleadings
%identified by the Court of Federal Claims in order to conclude that the court
%correctly dismissed Kam-Almaz's claim for breach of an implied-in-fact bailment
%contract.
%
%Finally, Kam-Almaz contends that the Court of Federal Claims ``treated the issue
%of the bailment as if such a claim were legally not viable,'' and in so doing
%disregarded the Supreme Court's opinions of \textit{Kosak v. United States}, 465
%U.S. 848, 104 S.Ct. 1519, 79 L.Ed.2d 860 (1984) and \textit{Hatzlachh Supply Co.
%v. United States}, 444 U.S. 460, 100 S.Ct. 647, 62 L.Ed.2d 614 (1980) (per
%curiam), as well as our court's decision in \textit{Acadia Technology, Inc. v.
%United States}, 458 F.3d 1327 (Fed.Cir. 2006). Br. Pl.-Appellant Kam-Almaz at
%14. Those opinions, however, merely suggested that the possibility of alleging
%an implied-in-fact bailment contract is not foreclosed when Customs detains
%property. In \textit{Kosak}, a case involving a claim under the Federal Tort
%Claims Act, the Court noted the prospect of other remedies, including the
%possibility of bringing suit against an individual Customs official who
%negligently damaged detained goods and the possibility of bringing suit ``[i]f
%the owner of property detained by the Customs Service were able to establish the
%existence of an implied-in-fact contract of bailment between himself and the
%Service....'' 465 U.S. at 860 \& n. 22, 104 S.Ct. 1519. In \textit{Hatzlachh},
%which similarly involved the effect of a statutory exception to the Federal Tort
%Claims Act, the Court observed that ``[t]he absence of Government tort liability
%has not been thought to bar contractual remedies on implied-in-fact
%contracts....'' 444 U.S. at 465, 100 S.Ct. 647. And finally, in \textit{Acadia},
%a case involving a takings claim under the Fifth Amendment, we noted ``that an
%owner [of property detained by Customs] might be able to bring a suit under the
%Tucker Act for money damages under a theory of breach of an implied-in-fact
%contract of bailment between the owner and Customs,'' but ``we express[ed] no
%opinion as to whether the facts [alleged] might support such a claim.'' 458 F.3d
%at 1334 n. 2. In the present case, it is unnecessary to opine whether under some
%hypothetical set of alleged facts an implied-in-fact bailment contract could
%conceivably arise from the detainment of property by Customs. We merely hold, on
%the specific facts alleged and for the reasons stated, that Kam-Almaz's
%complaint fails to state a claim for an implied-in-fact bailment contract.
%
%
%
%\readinghead{II}
%
%We next turn to Kam-Almaz's claim for a compensable taking. The Fifth Amendment
%to the United States Constitution provides in part, ``nor shall private property
%be taken for public use, without just compensation.'' The purpose of the Takings
%Clause is ``to bar Government from forcing some people alone to bear public
%burdens which, in all fairness and justice, should be borne by the public as a
%whole.'' \textit{Armstrong v. United States}, 364 U.S. 40, 49, 80 S.Ct. 1563, 4
%L.Ed.2d 1554 (1960).
%
%Kam-Almaz asserts that the Court of Federal Claims erred by dismissing his
%takings claim for failure to state a claim under RCFC 12(b)(6). According to
%Kam-Almaz, the government's seizure of his laptop constituted a physical taking
%for public use, for which just compensation is due. Kam-Almaz further contends
%that his claim alleges a compensable taking because, unlike cases such as
%\textit{Bennis v. Michigan}, 516 U.S. 442, 116 S.Ct. 994, 134 L.Ed.2d 68 (1996),
%a crime was not committed using the seized property. In addition, Kam-Almaz
%argues that the reasoning of the Court of Federal Claims was flawed because it
%did not define the extent of the government's ``police power.''
%
% The government, in response, contends that the Court of Federal Claims lacked
%jurisdiction over Kam-Almaz's takings claim because Kam-Almaz asserted that the
%seizure of his laptop was unlawful and unjust. Assuming the court had
%jurisdiction, the government asserts that the court correctly dismissed
%Kam-Almaz's claim. According to the government, the complaint failed to allege a
%constitutional taking because an authorized seizure of property pursuant to the
%government's police power is not for public use, and thus cannot be a taking. In
%addition, the government argues that property detained or seized by ICE
%officials is necessarily taken pursuant to the government's police power.
%
%As an initial matter, we agree with Kam-Almaz that the Court of Federal Claims
%possessed jurisdiction over his takings claim. The assertion in Kam-Almaz's
%complaint that he suffered ``an unjust and unlawful taking of his property,''
%Compl. at ¶ 22, could be read in at least two ways. On the one hand, as the
%government contends, it could be an assertion that the government's seizure was
%unauthorized, in which case the district court, not the Court of Federal Claims,
%would possess jurisdiction. \textit{See} \textit{Acadia}, 458 F.3d at 1331
%(``[The] Tucker Act does not create jurisdiction in the Court of Federal Claims
%for a party contesting the propriety of a seizure.''). On the other hand, as
%Kam-Almaz asserts, the complaint may assert that the government's seizure,
%although authorized, was compensable under the Fifth Amendment. Under the
%liberal standards applied to motions to dismiss, we conclude that the complaint
%does not dispute the propriety of the seizure itself, but rather seeks to
%challenge the failure of the government to compensate Kam-Almaz for a
%presumptively lawful seizure. Thus, the Court of Federal Claims properly
%exercised its jurisdiction over Kam-Almaz's complaint.
%
%Turning to the matter of the court's dismissal, however, we agree with the
%government that Kam-Almaz failed to state a Fifth Amendment takings claim. Our
%precedent is clear: ``Property seized and retained pursuant to the police power
%is not taken for a `public use' in the context of the Takings Clause.''
%\textit{AmeriSource Corp. v. United States}, 525 F.3d 1149, 1153 (Fed.Cir.2008);
%\textit{see also} \textit{Acadia}, 458 F.3d at 1331-32. Kam-Almaz appears to
%recognize as much, arguing that, ``[t]o the extent the Court considers
%\textit{Acadia} and \textit{AmeriSource} to be controlling, Kam-Almaz
%respectfully submits that those cases should be overruled en banc.'' Br.
%Pl.-Appellant Kam-Almaz at 43. Of course, a panel of this court is powerless to
%overrule a precedent. \textit{See} \textit{El-Shifa Pharm. Indus. Co. v. United
%States}, 378 F.3d 1346, 1352 (Fed. Cir.2004).
%
%In any event, Supreme Court precedent also supports the dismissal of Kam-Almaz's
%complaint. In \textit{Bennis}, the Court held that Mrs. Bennis, an innocent wife
%who had a property interest in a car that was forfeited after her husband used
%it in the commission of a crime, did not have a claim for a compensable taking.
%516 U.S. at 452-53, 116 S.Ct. 994. Mrs. Bennis's innocence did not factor into
%the Court's takings analysis. Rather, the Court held that ``[t]he government may
%not be required to compensate an owner for property which it has already
%lawfully acquired under the exercise of governmental authority other than the
%power of eminent domain.'' \textit{Id.} at 452, 116 S.Ct. 994. In
%\textit{AmeriSource}, we summarized the Court's holding in \textit{Bennis} as
%follows: ``[in] case[s] involving governmental seizure of property for law
%enforcement purposes, the [takings] inquiry remains focused on the character of
%the government action, not the culpability or innocence of the property
%holder.'' 525 F.3d at 1154. Accordingly,  Kam-Almaz's innocence does not convert
%ICE's seizure into a compensable taking under the Fifth
%Amendment.\readingfootnote{1}{Kam-Almaz relies on \textit{Shelden v. United
%States}, 7 F.3d 1022 (Fed.Cir.1993), in asserting that a takings claim may lie
%when property is seized by the government from an innocent owner pursuant to a
%criminal investigation. \textit{Shelden}, however, was decided before
%\textit{Bennis.} As we noted in \textit{AmeriSource}, ``[t]o the extent that
%[\textit{Shelden}] purports to create any rules with respect to innocent owners
%in the takings context, it plainly lacks force.'' 525 F.3d at 1156. Kam-Almaz's
%reliance on \textit{Shelden} is therefore not persuasive.}
%
%Kam-Almaz further asserts that the government's seizure was not an exercise of
%its police power. We disagree. Customs officers unquestionably have the
%authority to search and seize property at our nation's
%borders.\readingfootnote{2}{The parties do not dispute that Agent Muldowan of
%ICE served as a Customs officer in performing the acts alleged in Kam-Almaz's
%complaint. \textit{See} 19 C.F.R. \S\S~162.6, 162.21; \textit{see also} 19
%U.S.C. \S~1401(i). We therefore have no occasion in this case to distinguish
%between Customs and ICE concerning the authority of their officers to perform
%border searches and seizures.} \textit{See, e.g.}, 19 C.F.R. \S\S~162.6, 162.21.
%Under the relevant regulation, ``[p]roperty may be seized ... by any Customs
%officer who has reasonable cause to believe that any law or regulation enforced
%by Customs and Border Protection or Immigration and Customs Enforcement has been
%violated....'' \textit{Id.} \S~162.21. Lawful seizures performed pursuant to
%such authority necessarily fall within the government's power to police the
%border. The ``police powers ... are nothing more or less than the powers of
%government inherent in every sovereignty to the extent of its dominions.''
%\textit{The License Cases}, 46 U.S. 504, 583, 5 How. 504, 12 L.Ed. 256 (1847)
%(opinion of Taney, C.J.). As we have noted, ``[a]lthough the precise contours of
%the principle are difficult to discern, it is clear that the police power
%encompasses the government's ability to seize and retain property to be used as
%evidence in a criminal prosecution.'' \textit{AmeriSource}, 525 F.3d at 1153.
%
%We therefore hold that the Court of Federal Claims correctly dismissed
%Kam-Almaz's takings claim under RCFC 12(b)(6). Whatever claim Kam-Almaz may have
%against the United States, if any, it is not under a breach of an
%implied-in-fact contract or a takings theory.



%\readinghead{CONCLUSION}
%
%We have considered Kam-Almaz's remaining arguments and find them unpersuasive.
%The judgment of the Court of Federal Claims dismissing Kam-Almaz's complaint is
%affirmed.
%
%AFFIRMED
%
%\vskip\baselineskip

\opinion \textsc{Newman}, Circuit Judge, dissenting.\ldots

%Mr. Kam-Almaz, a United States citizen, returning to Dulles International
%Airport from international business travel, was accused of no wrongdoing, and
%his property was seized, detained, and damaged. He was without fault and without
%accusation. The court now condones this action and injury, denying all remedy
%for the destruction of his hard drive and the loss of valuable business
%information on his computer.
%
%Mr. Kam-Almaz was travelling abroad for his work with the International
%Institute for Psychosocial Development, a non-profit organization that
%specializes in post-disaster mental health education and consulting. On his
%return to Dulles Airport, upon routine Customs inspection his laptop computer
%was seized by a Customs agent, the agent stating that the computer files  would
%be reviewed because he was a ``person of interest.'' Mr. Kam-Almaz objected that
%the computer contained valuable business data, and the agent allowed him to copy
%one file but not all of the computer files. The agent told Mr. Kam-Almaz that
%the laptop would be returned within seven days, and he was given a form receipt
%that said his property would be returned within thirty days. The laptop was not
%returned for more than two months, after persistent inquiry by Mr. Kam-Almaz.
%When returned, the hard drive and many of the computer files had been destroyed.
%
%Mr. Kam-Almaz was accused of no wrongdoing, but his property was destroyed while
%in the custody of the government. The issue is not whether the government could
%have seized and detained his computer under its police power; the issue is
%whether the government, having taken possession of the computer and destroyed
%its hard drive and files, had any responsibility for the property and has any
%liability for the loss incurred. The court now condones and ratifies this injury
%to Mr. Kam-Almaz's property, and holds that he is without remedy. I respectfully
%dissent.
%
%
%
%\readinghead{DISCUSSION}
%
%This suit is brought under the Tucker Act. The government's position is that the
%government action is not a Fifth Amendment taking of property because the
%property was taken in exercise of police power. The government also states its
%action did not create an implied-in-fact bailment contract because the computer
%was seized without mutual consent, and that all contracts require consent. My
%colleagues agree, and since the Supreme Court has held that no remedy is
%available on a theory of tort, transfer to a district court is not available,
%depriving Mr. Kam-Almaz of all access to remedy. However, as I shall explain,
%the court errs in holding that Mr. Kam-Almaz is without Tucker Act remedy for
%the destruction of his property by the government.
%
%
%
%\readinghead{I. The Contract Claim}

The court holds that Mr. Kam-Almaz's claim for breach of an implied contract of
bailment cannot lie, because his complaint states that his property was seized.
The court reasons that a ``seizure, essentially by definition, lacks mutual
intent,'' and thus that the government cannot be liable on
contract theory because it took the property without consent. This holding is
contrary to law and precedent.

In several rulings, typified by \textit{Kosak v. United States}, 465 U.S. 848
(1984), the Court held, under circumstances of
Customs seizure, that remedy may be available on a theory of implied-in-fact
contract of bailment. In \textit{Kosak} the Customs Service seized antiques and
objects of art belonging to the plaintiff. The property was eventually returned,
but was damaged while in the custody of the Customs Service. The plaintiff
brought suit under the Federal Tort Claims Act. The Court held that the Tort
Claims Act did not apply to injury to property detained by the Customs Service,
but\ldots:
% that ``there exists at least one other remedial system,''
% including suit
%under the Tucker Act based on implied-in-fact contract:
\begin{quote}
[T]here exists at least one other remedial system that might
enable someone in petitioner's position to obtain compensation from the
Government. If the owner of property detained by the Customs Service were able
to establish the existence of an implied-in-fact contract of bailment between
himself and the Service, he could bring suit\ldots.
% under the Tucker Act, 28 U.S.C.
%\S~1491. \textit{See} \textit{Hatzlachh Supply Co. v. United States}, 444 U.S.
%460 [100 S.Ct. 647, 62 L.Ed.2d 614] (1980).
\end{quote}
The Court observed that the Kosak
property was ``seized,'' and did not
distinguish seizure from involuntary detention.

In contrast, my colleagues hold that since Mr. Kam-Almaz alleged a seizure of
his property, he is precluded from seeking relief under a contract theory. This
holding conflicts with \textit{Kosak}, and with \textit{Hatzlachh Supply Co.,
Inc. v. United States}, 444 U.S. 460 (1980),
where the Court held that ``the United States may be held liable for breach of
an implied contract of bailment when goods are lost while held by the United
States Customs Service (USCS) following their seizure for customs violations.''
The Court did not hold that ``seizure''
eliminated liability for breach of an implied contract of bailment of the seized
goods. My colleagues have departed from clear precedent.

[This is only the beginning of Judge Newman's extensive dissent.]

%These are not new theories of government obligation and liability. In
%\textit{Alliance Assurance Co. v. United States}, 252 F.2d 529 (2d Cir.1958),
%suit was brought in the district court under the Tucker Act for the value of
%goods which, ``while being inspected for entry into this country, disappeared
%from the possession of the United States Customs.'' \textit{Id.} at 531. The
%government was held liable for breach of an implied contract, the court
%explaining:
%
%\begin{quote} The obligation of the government was not artificially created by
%law but rather stemmed from an implied promise to redeliver the goods as soon as
%customs had checked them against the invoice. Such a promise need not be
%formalized in a written agreement or even made the subject of a specific
%conversation. It arises from the implied promise to return the goods to the
%lawful owner after the customs inspection has been completed. \end{quote}
%\textit{Id.} at 532. The court held that the government ``voluntarily undertook
%a bailment of the goods in question, a promise on its part to use due care
%during the term of the bailment can and should be implied.'' \textit{Id.}
%
%Contrary to precedent, my colleagues hold that ``because Kam-Almaz did not
%voluntarily deliver his property to the government, his complaint fails to
%allege any valid consideration.'' Maj. op. at 1369. In \textit{Alliance} the
%court explained that ``the owner's trusting him with the goods is a sufficient
%consideration to oblige him to a careful management,'' and that ``compelling
%reason to find consideration exists here because the bailment, although
%gratuitous, was compulsory and for the exclusive benefit of the bailee.''
%\textit{Id.} at 533. This reasoning applies in Customs' detention of Mr.
%Kam-Almaz's computer and files, for the bailment was compulsory and solely for
%the benefit of the government.
%
%The panel majority is incorrect in holding that Mr. Kam-Almaz's claim is barred
%because he ``did not voluntarily `deliver' his equipment to Muldowan,'' maj. op.
%at 1369, for, as the Supreme Court and other courts have explained, the
%involuntary detention with a promise to return the property confirms the
%implied-in-fact bailment. Mr. Kam-Almaz was given a written receipt for the
%laptop, stating in writing that it would be returned within thirty days. This is
%written confirmation of the contract of bailment.
%
%The fact that Mr. Kam-Almaz did not voluntarily part with his laptop does not
%absolve the government from liability for the injury to his property.
%Acquiescence in Customs' seizure for inspection of the property does not include
%authorization to destroy the property. As the Court stated in \textit{Kosak},
%the ``owner of property detained by the Customs Service'' has ``at least'' an
%implied-in-fact contract remedy.
%
% Ignoring this direct precedent, my colleagues dismiss Mr. Kam-Almaz's contract
%claim, seeking support in seizures of contraband or illegal activity. In
%\textit{Llamera v. United States}, 15 Cl.Ct. 593 (1988), the Court of Federal
%Claims held that the government was not liable on a bailment theory for loss of
%a vessel that was seized by the Coast Guard for violation of law. The Court of
%Federal Claims found that the Coast Guard had told the plaintiff that the Coast
%Guard ``was assuming no responsibility whatsoever for the vessel'' it had
%seized, and ``told plaintiff's brother that he could stay with the vessel or
%have someone board the vessel to watch it for him.'' 15 Cl.Ct. at 597-98. Such
%explicit statement of absence of responsibility for the vessel is far removed
%from the facts of this case.
%
%The panel majority also places inappropriate reliance on the criminal seizure in
%\textit{Alde, S.A. v. United States}, 28 Fed.Cl. 26 (1993), where the government
%seized an aircraft from the Dominican Republic, on its arrival in Puerto Rico,
%on warrant for suspected violation of law. While in government control, the
%aircraft was damaged in a hurricane. The court held that the government was not
%liable on either contract or takings theory, because:
%
%\begin{quote} The Customs Service eventually did seek forfeiture of the
%aircraft. In light of these facts, it is difficult to fathom how plaintiff
%believed the Customs Service was indicating that it would safeguard the aircraft
%and its contents for later return to plaintiff. Quite to the contrary, the
%Customs Service was actively seeking to permanently deprive plaintiff of
%ownership of the aircraft through forfeiture proceedings. \end{quote} 28 Fed.Cl.
%at 32. The \textit{Alde} ruling provides no support for this court's holding
%that the government has no liability for the destruction of Mr. Kam-Almaz's
%detained property, with no accusation of wrongdoing, before or after the
%detention.
%
%The record supports Mr. Kam-Almaz's position that the Customs agent stated the
%intent to return his property in a short time. The Customs agent was told that
%the computer held important and valuable business information, for the
%government permitted Mr. Kam-Almaz to copy one but not all of his files. The
%Complaint states:
%
%\begin{quote} At the time of seizure of his computer, Agent Craig Moldowan
%orally promised to return it to the Plaintiff within seven (7) days after a
%security review of its contents was performed by the Defendants. Agent Moldowan
%signed a written Customs Form 6051D, evidencing his written offer not to detain
%the Plaintiff's computer for longer than thirty days.... The Defendants breached
%their implied contract with Plaintiff when they failed to use due care in a
%prompt manner in the handling of the Plaintiff's computer, due to a Government,
%operator-caused system crash.... In addition, the Defendants failed to mitigate
%these damages by simply allowing the Plaintiff to make a copy of his operating
%software and data files, as he requested the defendants do. \end{quote} Applying
%law and precedent to the pleadings, Mr. Kam-Almaz has stated a claim for relief
%on an implied-in-fact contract theory. The complaint was improperly dismissed
%for failure to state a claim.
%
%
%
%\readinghead{II. The Takings Claim}
%
%The panel majority also dismisses Mr. Kam-Almaz's Fifth Amendment takings claim,
%on the ground that his property was ``seized and retained pursuant to the police
%power.'' Maj. op. at 1371. However, this court has recognized that ``it is
%insufficient to avoid the burdens imposed by the Takings Clause simply to invoke
%the `police powers' of the state, regardless of the respective benefits to the
%public and burdens on the property owner.'' \textit{Acadia Tech., Inc., Global
%Win Tech., Ltd. v. United}  \textit{States}, 458 F.3d 1327, 1330, 1332-33 (Fed.
%Cir.2006) (``[A] taking does not result simply because the government acted
%unlawfully, nor does a takings claim fail simply because the government's
%conduct is subject to challenge as unlawful.'').
%
%Although protection of the nation's borders is a police activity, when the
%government in its performance injures an innocent person, that person is not
%required to ``bear public burdens which, in all fairness and justice, should be
%borne by the public as a whole.'' \textit{Armstrong v. United States}, 364 U.S.
%40, 49, 80 S.Ct. 1563, 4 L.Ed.2d 1554 (1960). Mr. Kam-Almaz is not excluded from
%access to his Fifth Amendment rights, for the Fifth Amendment applies to
%government actions which seize and injure private property. \textit{See, e.g.},
%\textit{Kimball Laundry Co. v. United States}, 338 U.S. 1, 3-4, 69 S.Ct. 1434,
%93 L.Ed. 1765 (1949) (ordering just compensation when the government took
%temporary possession of a laundry plant); \textit{United States v. General
%Motors Corp.}, 323 U.S. 373, 375, 65 S.Ct. 357, 89 L.Ed. 311 (1945) (requiring
%compensation from the government for using a portion of a building leased by the
%plaintiff).
%
%The cases relied on by the panel majority all involve a seizure of contraband or
%other violation of law. Neither Mr. Kam-Almaz nor his property has been accused
%of or implicated in any wrongdoing. In \textit{Acadia, supra}\textit{}, the
%seized property allegedly bore a counterfeit trademark. In \textit{AmeriSource
%Corp. v. United States}, 525 F.3d 1149, 1150 (Fed.Cir.2008), drugs were seized
%pursuant to charges of ``conspiracy, unlawful distribution of prescription
%pharmaceuticals, operating an unregistered drug facility, and conspiracy to
%commit money laundering.'' In \textit{Bennis v. Michigan}, 516 U.S. 442, 116
%S.Ct. 994, 134 L.Ed.2d 68 (1996), the Court upheld the forfeiture of an
%automobile that had been used in criminal activity, although the other owner of
%the automobile was innocent. These cases all relate to confiscation or seizure
%based on unlawful acts. In contrast, no wrong by Mr. Kam-Almaz or his computer
%was alleged by the government.
%
%It is incorrect to apply criminal law when no crime is alleged. Precedent
%supports Mr. Kam-Almaz's pleading of a takings claim.
%
%
%
%\readinghead{III. Other Possible Remedies?}
%
%The panel majority refers to ``whatever claim Kam-Almaz may have against the
%United States,'' as if he simply chose the wrong forum. However, Customs injury
%cannot be remedied by tort claim, for in \textit{Kosak} the Court established
%that ``the Tort Claims Act does not cover suits alleging that customs officials
%injured property that had been detained by the Customs Service.'' 465 U.S. at
%862, 104 S.Ct. 1519. As mentioned \textit{supra}, the \textit{Kosak} Court
%stated that ``at least'' the Tucker Act might provide a remedy. \textit{Id.} at
%861 n. 22, 104 S.Ct. 1519.
%
%Whatever the remedial theory, the government is not absolved of responsibility
%for its misfeasance in dealing with its citizens. The Court of Claims ``holds
%and speaks a nation's conscience,'' \textit{see} tribute to Chief Judge Peele,
%48 Ct. Cl. XXV (Feb 11, 1913). The court has fulfilled this proud tradition
%since its inception:
%
%\begin{quote} A unique and permanent contribution that the Court of Claims has
%made over the span of its long life as a public institution is in how it helps
%make Government officials accountable to the citizens whose servants they are,
%but whose relationship to their masters is sometimes forgotten. In helping to
%inspire a high standard of conduct for Government officials, it serves the
%nation well. If there is a constant thread running through the court's
%decisions, it would seem to be in holding the Government  and its officials to a
%strict code of conduct in their relations with citizens. \end{quote} Judge
%Marion Bennett, in \textit{The United States Court of Claims: A History} 170-72
%(1978). The History explained that ``[s]uch a court is the flower of a free
%society.'' \textit{Id.}

Persons injured by unjustified damage to their property detained by Customs
agents are not excluded from access to the courts. This court's endorsement of
such exclusion misconstrues traditional and constitutional theory, and denies
the responsibility carved into this court's entrance, that: ``It is as much a
duty of government to render prompt justice against itself, in favor of
citizens, as it is to administer the same, between private individuals.''
Abraham Lincoln, Dec. 3, 1861.

I respectfully dissent.

